Attribute-classification accuracy can misrepresent the identification value of a soft biometric. We present a controlled framework separating oracle gains, prediction errors, evidence dependence and writer-associated information, and apply it to BanglaHand: 42 writers, six representations and 120 identification settings. Normalization uses selected gallery references only; query-handedness models exclude the queried writer from training and calibration. Known handedness improves top-1 accuracy by 6.1 percentage points, compared with 6.9 points for an expected known random partition. With query handedness predicted at 65.5\% accuracy on self-chosen paragraphs, calibrated fixed-weight fusion changes identification accuracy by -0.6 points, properly nested learned weighting by -0.08, and hard filtering by -14.3. Neither fusion rule establishes a gain in the designated comparisons. Filtering and margin diagnostics explain how joint correctness and confidence affect ranking. Random writer labels expose writer-associated information in enrolled-writer attribute evaluation. An external CSAFE evaluation uses 96 writers and separate gallery/query capture sessions. Handedness accuracy is 67.7\% for same-prompt and 58.3\% for cross-prompt session-2 queries; none of eight primary fusion comparisons establishes a Holm-adjusted gain. Session-3 and gallery-subsample analyses provide secondary checks. The framework provides reusable controls for evaluating auxiliary attributes in biometric candidate ranking, helping researchers distinguish oracle opportunity from operational benefit and make fusion decisions using measured identification utility.
